\subsection{Factoring Quadratic Trinomials: the AC Method} 
When we multiply two linear binomials (in one variable or homogeneous in the same two variables), the result is a quadratic trinomial.
\begin{example}
Multiply $(2x-3)(x+2)$, using the rectangle method and the FOIL shortcut.
\begin{lpic}{}{4}
\end{lpic}
\end{example}
\noindent When we multiply, we must \makeblank. When we factor, we must reverse this, i.e. ``\makeblank[1in]'' the middle terms.

Observing the multiplication, we see that the middle terms (to be combined) have the same \makeblank[1in] as the first and last terms. This gives us the following technique.
\begin{procedure}{The AC Method}
To factor a quadratic trinomial:
\begin{enumerate*}
\item Multiply the coefficients of the first and last terms.
\item Find two numbers that have the same product, and whose sum is the coefficient of the middle term.
\item Rewrite the middle term as a sum of two terms using these coefficients.
\item The polynomial now has four terms. Factor by grouping.
\end{enumerate*}
If we write a quadratic trinomial as $ax^2+bx+c$ (or $ax^2+bxy+cy^2$), the two numbers we seek in step 2 have product $ac$ and sum $b$.
\end{procedure}
\begin{example}
Factor $2x^2+x-6$ using the AC Method.
\begin{lpic}{}{4}
\end{lpic}
\end{example}
\begin{example}
Factor $6x^2-11xy+3y^2$ using the AC Method.
\begin{lpic}{}{4}
\end{lpic}
\end{example}
